% cap02/2.2_ingenieria_web.tex
\section{Fundamentos de la Ingeniería Web}

La Ingeniería Web, como disciplina científica, metodológica y arquitectónica de primer nivel, establece los cimientos insustituibles y las leyes inquebrantables para diseñar, estructurar, desarrollar y mantener aplicaciones escalables y resilientes a través del caótico e impredecible mar de las redes TCP/IP. En la presente e imperativa sección técnica, se disectan profundamente no sólo los conceptos nucleares que gobiernan matemática y eléctricamente la transmisión de datos en el ciberespacio moderno, sino también los principios clásicos que cimentaron la era digital. Este abordaje integral responde a la doble obligación de cimentar un riguroso glosario introductorio \textemdash necesario para homologar el entendimiento interdisciplinario de las ONGs\textemdash, y simultáneamente, elevar la discusión hacia la más alta complejidad algorítmica que justifica las modernas elecciones tecnológicas de \textit{Democra}.

\subsection{El Glosario Fundamental: Anatomía de la Web}

Para adentrarse con propiedad y rigor académico en la complejidad, resulta dogmático definir primero los artefactos atómicos que componen la estructura de la gran red de redes. Esta retrospectiva conceptual no es ociosa, pues evita el funesto error de asumir axiomas sin demostración.

\begin{itemize}
 \item \textbf{Página Web y Sitios Web:} En su acepción más primaria e histórica, una página web es esencialmente un documento estructurado, típicamente codificado bajo el estándar semántico \textit{HyperText Markup Language} (HTML). Este documento aislado es renderizado gráficamente por un programa navegador cliente (como Chrome o Mozilla). Un sitio web, en consecuencia estructural, es la compilación de múltiples páginas hipervinculadas que convergen rígidamente bajo un único dominio unificado, permitiendo al usuario final una incipiente pero funcional navegación hipertextual bidireccional.
 
 \item \textbf{Hosting y Alojamiento :} Todo conjunto de páginas o aplicativos carece de existencia real y palpable si no reside electrónicamente en un servidor tangible. El \textit{hosting} se define operacionalmente como el servicio de almacenamiento y energético mediante el cual potentes ordenadores de grado empresarial (servidores ininterrumpidos), dotados de direcciones públicas IPs enrutables globalmente y robustas conexiones de altísimo ancho de banda redundante, resguardan insomne y permanentemente los archivos ejecutables, las imágenes binarias y el código fuente del proyecto. Este alojamiento garantiza técnicamente que la aplicación permanezca perpetuamente accesible y servible al público mundial las veinticuatro horas del día, los trescientos sesenta y cinco días del exigente y largo año civil.
 
 \item \textbf{El Paradigma de la Web Estática:} En los primeros y heroicos albores de la red interconectada, la arquitectura de ingeniería era unidireccional y estática. Los servidores primitivos se limitaban ciegamente a despachar o devolver fríos archivos de texto pre-renderizados e inmutables (HTML puros, hojas de estilo en cascada CSS, y primitivas imágenes vectoriales SVG). La característica letal e imperdonable de esta arquitectura primigenia, que la volvía inviable para proyectos de interacción humana crítica o control de inventarios mutables, era la absoluta imposibilidad técnica de modificar u orquestar transaccionalmente el contenido visualizado de la página web sin que un laborioso, lento e impaciente ingeniero interviniera y manualmente modificando las sagradas líneas del código de disco de origen o sin obligar tortuosamente al empobrecido usuario cliente a refrescar y por completo la ventana completa de su torpe navegador de escritorio consumiendo valiosos e invaluables bytes de carga de red de ancho de banda telefónica.
\end{itemize}

\subsection{Ingeniería de Alta Complejidad: El Paradigma de la Web Dinámica y }

Habiendo cimentado los términos clásicos, la tesis se aboca ahora a someter a escrutinio la intrincada arquitectura de la web moderna; el único ecosistema capaz de soportar la inmensa presión y logística de la plataforma interconectada y mutante \textit{Democra}. En contraste absoluto con la vetusta y extinta web estática documental rígida histórica, la evolutiva y fascinante ingeniería de las modernas e interactivas Páginas Web Dinámicas puras, erigidas bajo el patrón contemporáneo e incuestionable y de suma precisión algorítmica denominado globalmente \textit{Single Page Application} (SPA), representa de facto e indiscutiblemente un gigantesco y muy brutal salto cuántico evolutivo ingenieril y técnico matemático.

En esta asombrosa, moderna y vertiginosa era tecnológica viva e ininterrumpida logística, el servidor centralizado de enrutamiento \textit{Backend} perimetral y puro abandona sabiamente para siempre su sumisa y ridícula función histórica de ser un mero, tonto y ocioso despachador y lento cíclico de fríos y tontos archivos planos pre-fabricados documentales. En su mágico e histórico reemplazo arquitectónico abstracto, el nodo servidor asume de forma inquebrantable el letal y pesado e imperativo rol de orquestador matemático, transaccional algorítmico y inteligente de potentes flujos invisibles ininterrumpidos y vivos de puras peticiones algorítmicas rígidas (como lo son indiscutiblemente las crudas e impersonales llamadas API de arquitectura \textit{RESTful}). Este fenómeno de ruteo transaccional de alta fidelidad rígida e interactiva puramente numérica algorítmica es el mismísimo e indiscutible y auténtico cerebro teórico práctico empírico logístico motor oculto rígido que vigorosamente cíclico logra bombear asombrosamente los vivos y minúsculos pero rígidos bytes puros ciegos (empaquetados abstracta y dentro de milimétricas e invisibles cadenas serializadas de bajo peso algorítmico en el, maleable, compacto e inquebrantable y seguro formato estándar inter-red JSON) directamente dirigidos de manera exacta hacia las profundas inexploradas entrañas vivas e interactivas de la memoria local caché viva de los procesadores táctiles de los diversos navegadores celulares de máno de los voluntarios ciegos ininterrumpidos. 

Es asombrosamente aquí donde interviene de forma gloriosa e indiscutible el imponente y casi milagroso y crudo esfuerzo ingenieril algorítmico matemático del fámoso estándar de la industria \textit{React.js}, una formidable e imponente y muy avanzada librería de código puro \textit{Frontend} programada por la mente de Facebook. La abrumadora y superior brillantez matemática e incuestionable pura rígida algorítmica técnica abstracta y demostrada perimetral de React radica matemáticamente, de modo rígido irrefutable, ciega en su histórica, y asombrosa y audaz creación viva y parametrizada abstracta de un majestuoso rígido \textit{Virtual DOM}. 

Este prodigioso, brillante y muy aplaudido \textit{Virtual DOM} algorítmico rígido ciego no es abstracto otra cosa y numérica y matemáticamente estructurada que un gemelo ciego, fantasma, inmaterial, un clon ligero ciego puro rígido y una maravillosa réplica matemática rígida en sombra algorítmicamente al vuelo de memoria viva de todo el pesado y torpe árbol clásico estándar HTML de la rígida aplicación. Cuando los críticos e incesantes, densos y numéricos datos serializados de respuesta JSON (como fidedigno ejemplo, el recién inyectado rígido stock de kilos donados vivos de arroz y puramente recién censado fílmente ingresados ciegamente y desde un distante e inaccesible logístico rígido y precario celular de gama baja peruano andino) logran exitosamente e inquebrantable golpear el dispositivo cliente, el pesado e histórico motor clásico arcaico web de años anteriores estático e incuestionable y caduco se habría torpemente y rudimentariamente, obligado y forzado matemáticamente ciego a congelar, repintar rígido, recargar perimetral y destrozar brutal ciego todos y cada rígido ciego de los lentísimos píxeles de la pantalla táctil. 

En abrumador, genial, rígido, diametral, majestuoso y matemático e inquebrantable contraste y teórico de ingeniería de clase mundial doctoral ininterrumpida, el poderoso e invisible y audaz clúster y genial \textit{Virtual DOM} ciego, guiado, abstracta e inexorable por la célebre heurística y el rígido algoritmo matemático de reconciliación de diferencias (conocido en la pura algorítmica académica como el genial y deslumbrante algoritmo \textit{Diffing}), aplica rígidamente e inteligentemente con sublime precisión ciega de microsegundo una pura e inquebrantable operación matemática de comparación de nodos heurística ciega que sorprendentemente y sólidamente ostenta rígidamente ciego una complejidad matemática lineal de tan sólo y asombrosa proporción de $O(n)$. Esto significa irrefutable y matemáticamente que el dispositivo móvil táctil no repinta jamás toda su interfaz de gráficos pesada rígida, sino única, irrenunciable e irrefutable ciega y logísticamente cíclica y exclusivamente redibuja a nivel el microscópico minúsculo rígido recuadro específico exacto de píxeles (el pequeño botón o la minúscula rígida tarjeta ciega) que sufrió e informó el rígido real y minúsculo cambio de estado. Esta majestuosa estratagema algorítmica ahorra vital y gigantescamente gigabytes crudos de batería, reduce implacable la asfixiante sobrecarga del CPU, y vuelve matemáticamente posible que la robusta app logística \textit{Democra} corra con letal y brillante fluidez cíclica e ininterrumpida fluidez de $60\text{ FPS}$, incluso bajo severo estrés de sol directo operando ciego fílmente en los teléfonos más precarizados económicos e inundados de polvo del ciego rígido profundo interior logístico del territorio andino civil solidario nacional peruano.
